\documentclass{situationist-preprint}

\title{Virtuous and Vicious Circles in Transcendental Method}
\msid{PHE-18-b}
\msfamily{PHE-18}
\msdiscipline{Phenomenology}
\mstype{Companion manuscript}
\mskeywords{circularity; transcendental method; rigour; epistemic circularity; hermeneutic circle; Fink}

\begin{document}
\maketitle

\begin{abstract}
Transcendental philosophy investigates the conditions of knowledge with the knowledge whose conditions it investigates, and so moves in a circle. PHE-18-a shows that phenomenology's circle with respect to rigour is virtuous, by the criterion that the standard discriminates and that its self-application is a test that could fail. We apply the criterion to three circles of the tradition: Hegel's objection to critical epistemology, Heidegger's hermeneutic circle, and Fink's problem of the phenomenologizing onlooker. We show that the criterion separates them along the line its authors intended, and give a simple sufficient condition for viciousness: a circle is vicious when the defining discipline's satisfaction of the standard follows from the definition alone.
\end{abstract}

\section{The criterion}

Husserl's claim that philosophy can be a rigorous science \citep{Husserl1911} is defended in PHE-18-a by way of a circle. Definition~2.7 of PHE-18-a calls a discipline $p$ that asserts a definition $\delta$ of a property $R$ a \emph{virtuous circle} if (i) $p$ has $R$, (ii) some discipline lacks $R$, and (iii) if $\delta$ were not evident, $p$ would lack $R$. Theorem~3.4 shows that phenomenology with its definition of rigour meets the three conditions, and Proposition~3.5 exhibits a vicious alternative, rigour as agreement with phenomenology, which discriminates but fails (iii).

\begin{proposition}\label{p:suff}
If a discipline's having $R$ follows from $\delta$ together with the bare fact that the discipline asserts $\delta$, the circle fails condition (iii) and is vicious.
\end{proposition}

\begin{proof}
Suppose $p$'s having $R$ follows from $\delta$ and the fact that $p$ asserts $\delta$. Neither premise concerns the evidence of $\delta$. So $p$ has $R$ whether or not $\delta$ is evident, and the counterfactual of condition (iii) is false.
\end{proof}

Proposition~3.5 of PHE-18-a is an instance: that phenomenology asserts only what phenomenology asserts follows from nothing. Condition (iii) is thus the substantive requirement. It fails exactly when the self-application is decided by the form of the definition rather than by the evidence of what the discipline asserts.

\section{Three circles}

\paragraph{Hegel's objection.} Hegel compared the demand that cognition be examined before it is used to the resolve of learning to swim before venturing into the water \citep[§10]{Hegel1830}. The objection is that any examination of cognition is itself cognition, so the examination either presupposes what it examines or never begins. The criterion answers the objection without accepting its conclusion. An examination of cognition that uses cognition is circular, but the circle is virtuous if its standard discriminates and its self-application could fail. Husserl's definition of rigour is of this kind (Lemma~3.2 of PHE-18-a): it uses evidence to establish what evidence is, and phenomenology would fail it if reflection did not in fact present rigorous knowledge as evident and re-verifiable.

\paragraph{The hermeneutic circle.} Heidegger held that understanding moves within a fore-structure of anticipations, and that what is decisive is not to escape the circle but to enter it in the right way, by working out the fore-structure in terms of the things themselves \citep[§32]{Heidegger1927}. Read through the criterion, \enquote{the right way} is condition (iii): an interpretation whose anticipations can be disappointed by the text or the matter is in a virtuous circle, and one that reads its anticipations into whatever it meets is in a vicious one, of the kind described by Proposition~\ref{p:suff}.

\paragraph{The phenomenologizing onlooker.} Fink argued that a transcendental theory of method must reflect on the reflecting ego, whose activity is presupposed by every phenomenological description, and that this reflection cannot be performed in the natural attitude or justified from outside the reduction \citep{Fink1995}. The circle here is that the reduction is clarified only by a further reduction. It is virtuous by the criterion if the clarification can fail: if reflection on the onlooker could reveal that the descriptions it produced were not evident, the method would be shown deficient by its own means. Fink's insistence that the transcendental theory of method is a critique, and not a self-confirmation, is a statement of condition (iii).

\begin{objection}
Condition (iii) is a counterfactual about the evidence of the definition. If the definition is in fact evident, the counterfactual's antecedent is never realized, and the test is never run.
\end{objection}

\begin{reply}
A test need not fail to be a test. What matters is that the outcome depends on something other than the definition's form, so that the investigation could have come out otherwise. A measuring instrument calibrated against itself is useless only if its readings are fixed by the calibration; one whose readings are answerable to what it measures is not, even when it reads correctly. The same distinction underlies the acceptance of epistemic circularity in establishing the reliability of perception, where track-record arguments use perception and could, in principle, reveal its unreliability \citep{Alston1986}.
\end{reply}

\section{Conclusion}

The criterion of PHE-18-a is general. It sorts the circles of transcendental method into those whose self-application could fail and those decided by their form, and Proposition~\ref{p:suff} gives the mark of the latter. Phenomenology's circle, Heidegger's circle entered in the right way, and Fink's transcendental theory of method fall on the virtuous side.

\begin{thebibliography}{9}\small
\bibitem[Alston(1986)]{Alston1986} Alston, W.~P. (1986). Epistemic circularity. \emph{Philosophy and Phenomenological Research}, 47(1), 1--30.
\bibitem[Fink(1995)]{Fink1995} Fink, E. (1995). \emph{Sixth Cartesian Meditation: The Idea of a Transcendental Theory of Method} (R. Bruzina, Trans.). Indiana University Press. (Original work written 1932.)
\bibitem[Hegel(1830)]{Hegel1830} Hegel, G.~W.~F. (1830). \emph{Enzyklop\"adie der philosophischen Wissenschaften im Grundrisse}. 3rd ed. Heidelberg.
\bibitem[Heidegger(1927)]{Heidegger1927} Heidegger, M. (1927). \emph{Sein und Zeit}. Niemeyer.
\bibitem[Husserl(1911)]{Husserl1911} Husserl, E. (1911). Philosophie als strenge Wissenschaft. \emph{Logos}, 1, 289--341.
\end{thebibliography}

\end{document}
